\chapter{L'environnement de développement web}

\begin{synthese}[Au programme]
Comprendre l'architecture \textbf{client/serveur} et le cycle requête/réponse HTTP,
identifier le rôle d'un \textbf{serveur web} (Apache, Nginx, IIS), installer et
configurer un \textbf{environnement intégré} (WampServer, XAMPP, LAMP), repérer le
\textbf{répertoire racine} et le \textbf{port d'écoute}, puis \textbf{héberger un
site en local} avant d'évoquer la mise en ligne.
\end{synthese}

\section{L'architecture client/serveur}

Tout site web dynamique repose sur un dialogue entre deux acteurs : un
\textbf{client} et un \textbf{serveur}. Avant d'écrire la moindre ligne de PHP, il
faut bien comprendre qui fait quoi.

\begin{definition}
Dans le modèle \textbf{client/serveur}, le \textbf{client} demande un service et le
\textbf{serveur} le fournit. Sur le Web, le client est le \textbf{navigateur}
(Chrome, Firefox, Edge\dots) et le serveur est un \textbf{serveur web} qui héberge
les pages et répond aux demandes.
\end{definition}

\subsection{Le client : le navigateur}

Le navigateur est le logiciel installé sur la machine de l'utilisateur. Son rôle :
\begin{itemize}
  \item envoyer une \textbf{requête} lorsqu'on saisit une adresse (URL) ou qu'on
        clique sur un lien ;
  \item recevoir la \textbf{réponse} du serveur (du code HTML, des images, du CSS\dots) ;
  \item \textbf{afficher} la page à l'écran après l'avoir interprétée.
\end{itemize}
Le navigateur ne « voit » jamais le code PHP : il ne reçoit que du HTML déjà
fabriqué par le serveur.

\subsection{Le serveur web}

\begin{definition}
Un \textbf{serveur web} est un logiciel (et la machine qui l'héberge) qui reste en
permanence à l'écoute des requêtes des clients, recherche la ressource demandée
(une page, une image\dots) et la renvoie au navigateur.
\end{definition}

Les serveurs web les plus répandus sont :
\begin{itemize}
  \item \textbf{Apache} (\emph{Apache HTTP Server}) : libre, gratuit, le plus utilisé
        pour apprendre et pour de nombreux sites ;
  \item \textbf{Nginx} : réputé pour sa rapidité et sa faible consommation, très
        utilisé sur les gros sites ;
  \item \textbf{IIS} (\emph{Internet Information Services}) : le serveur web de
        Microsoft, intégré à Windows Server.
\end{itemize}

\begin{remarque}
Le mot « serveur » désigne tantôt le \emph{logiciel} (Apache), tantôt la
\emph{machine} (un ordinateur puissant allumé en permanence). Sur ton propre PC, le
même ordinateur peut jouer à la fois le rôle de client \emph{et} de serveur : c'est
exactement ce qu'on fait quand on installe WampServer pour tester un site en local.
\end{remarque}

\subsection{Le cycle requête / réponse HTTP}

Le navigateur et le serveur communiquent grâce à un langage commun : le protocole
\textbf{HTTP} (\emph{HyperText Transfer Protocol}). L'échange suit toujours le même
cycle en quatre temps.

\begin{enumerate}[label=\arabic*)]
  \item \textbf{Requête.} L'utilisateur saisit \code{http://www.onbuch.cm/index.php}.
        Le navigateur envoie une \textbf{requête HTTP} au serveur (« donne-moi la
        page \code{index.php} »).
  \item \textbf{Traitement.} Le serveur web reçoit la requête. Si la page contient
        du PHP, il fait \textbf{exécuter} le code par le préprocesseur PHP (qui peut
        interroger une base de données).
  \item \textbf{Réponse.} Le serveur renvoie une \textbf{réponse HTTP} contenant la
        page \textbf{HTML} finale (le résultat, plus le code PHP).
  \item \textbf{Affichage.} Le navigateur interprète le HTML reçu et affiche la page
        à l'utilisateur.
\end{enumerate}

\begin{center}
\small
\begin{tikzpicture}[font=\sffamily\small]
  % boîtes
  \node[draw=onbblue, very thick, rounded corners, fill=onbblue!6,
        minimum width=3.1cm, minimum height=1.5cm, align=center]
        (client) at (0,0) {\textbf{CLIENT}\\\footnotesize Navigateur\\\footnotesize (Chrome, Firefox)};
  \node[draw=onbo, very thick, rounded corners, fill=onbsoft,
        minimum width=3.1cm, minimum height=1.5cm, align=center]
        (serveur) at (8,0) {\textbf{SERVEUR WEB}\\\footnotesize Apache + PHP\\\footnotesize + base de données};
  % flèche requête
  \draw[->, >=latex, very thick, onbblue]
        ([yshift=4pt]client.east) -- node[above, midway, onbblue]{\footnotesize 1. Requête HTTP}
        node[below, midway, onbmuted]{\footnotesize \code{GET /index.php}} ([yshift=4pt]serveur.west);
  % flèche réponse
  \draw[->, >=latex, very thick, onbo]
        ([yshift=-10pt]serveur.west) -- node[above, midway, onbo]{\footnotesize 3. Réponse HTTP}
        node[below, midway, onbmuted]{\footnotesize page HTML} ([yshift=-10pt]client.east);
  % traitement
  \node[onbmuted, align=center, font=\sffamily\footnotesize] at (8,-1.6)
        {2. Exécution du PHP\\(génération du HTML)};
\end{tikzpicture}
\end{center}

\fig{Le cycle requête/réponse HTTP : le client demande, le serveur exécute puis répond.}

\section{Les environnements intégrés de développement web}

Pour développer un site PHP, il faut au minimum trois logiciels qui fonctionnent
ensemble. Plutôt que de les installer un par un, on utilise un \textbf{environnement
intégré} qui les regroupe en une seule installation.

\begin{definition}
Un \textbf{environnement intégré de développement web} est une suite logicielle qui
réunit, prêts à l'emploi, un \textbf{serveur web}, un \textbf{SGBD} (système de
gestion de base de données) et un \textbf{préprocesseur} de langage côté serveur.
\end{definition}

\subsection{Les trois composants}

Quel que soit l'environnement choisi, on retrouve toujours les mêmes briques :

\begin{center}\small
\begin{tabular}{@{}lll@{}}
\toprule
\textbf{Composant} & \textbf{Logiciel typique} & \textbf{Rôle}\\
\midrule
Serveur web    & Apache              & Reçoit les requêtes, renvoie les pages\\
SGBD           & MySQL / MariaDB     & Stocke et gère les données (base de données)\\
Préprocesseur  & PHP                 & Exécute le code et fabrique le HTML\\
\bottomrule
\end{tabular}
\end{center}

\begin{remarque}
On parle souvent d'\textbf{outil AMP} : \textbf{A}pache + \textbf{M}ySQL +
\textbf{P}HP. La première lettre indique le système d'exploitation : \textbf{W}AMP
(Windows), \textbf{L}AMP (Linux), \textbf{M}AMP (macOS).
\end{remarque}

\subsection{Les environnements les plus courants}

\begin{center}\small
\begin{tabular}{@{}llll@{}}
\toprule
\textbf{Environnement} & \textbf{Système} & \textbf{Répertoire racine} & \textbf{Particularité}\\
\midrule
WampServer & Windows        & \code{www}    & Très répandu dans les lycées\\
XAMPP      & Multi-plateforme & \code{htdocs} & Marche sous Windows, Linux, macOS\\
EasyPHP    & Windows        & \code{www}    & Léger, simple à installer\\
LAMP       & Linux          & \code{/var/www/html} & Installé paquet par paquet\\
\bottomrule
\end{tabular}
\end{center}

\begin{remarque}
\textbf{LAMP} n'est pas un logiciel unique à télécharger : c'est l'ensemble
\textbf{L}inux + \textbf{A}pache + \textbf{M}ySQL + \textbf{P}HP qu'on installe
généralement avec le gestionnaire de paquets (par exemple
\code{sudo apt install apache2 mysql-server php}). WampServer, XAMPP et EasyPHP, eux,
s'installent en un seul fichier d'installation.
\end{remarque}

\section{Installer et configurer l'environnement}

L'installation d'un environnement comme WampServer ou XAMPP se résume à télécharger
le programme depuis le site officiel, puis à suivre l'assistant. Une fois installé,
le logiciel place une icône dans la \textbf{zone de notification} (à côté de
l'horloge) qui permet de piloter les services.

\begin{methode}
\textbf{Installer WampServer et tester l'installation.}
\begin{enumerate}[label=\arabic*)]
  \item \textbf{Télécharger} WampServer depuis le site officiel
        (\code{wampserver.com}) en choisissant la version 32 ou 64 bits qui
        correspond à ton Windows.
  \item \textbf{Lancer l'installeur} et accepter les options par défaut (dossier
        d'installation \code{C:\textbackslash wamp64}). Laisser installer les
        composants Apache, MySQL et PHP.
  \item \textbf{Démarrer WampServer} : une icône apparaît près de l'horloge. Elle
        devient \textbf{verte} quand tous les services sont actifs (orange = un
        service manque, rouge = aucun service).
  \item \textbf{Tester} : ouvrir le navigateur et saisir \code{http://localhost}.
        La page d'accueil de WampServer doit s'afficher : l'installation est réussie.
  \item \textbf{Créer une première page} : déposer un fichier \code{essai.php} dans
        le dossier \code{www}, puis ouvrir \code{http://localhost/essai.php}.
\end{enumerate}
\end{methode}

\begin{attention}
Sur certains réseaux de lycées ou de cybercafés, un autre logiciel (Skype, un
antivirus, ou un autre serveur) peut déjà utiliser le port 80 et empêcher Apache de
démarrer (icône orange). Dans ce cas, il faut soit fermer ce logiciel, soit
\textbf{changer le port d'écoute} d'Apache (voir plus loin).
\end{attention}

\section{Démarrer et arrêter les services}

Un \textbf{service} est un logiciel qui tourne en arrière-plan. Dans WampServer ou
XAMPP, les deux services essentiels sont \textbf{Apache} (le serveur web) et
\textbf{MySQL} (le SGBD).

\begin{itemize}
  \item \textbf{WampServer} : un clic \emph{gauche} sur l'icône ouvre un menu avec
        \emph{Start All Services}, \emph{Stop All Services} et \emph{Restart All
        Services}. Le menu \emph{Apache} (et \emph{MySQL}) permet aussi de démarrer
        ou arrêter chaque service séparément.
  \item \textbf{XAMPP} : le \emph{Control Panel} affiche chaque module avec un bouton
        \emph{Start} / \emph{Stop}. Un voyant vert signale qu'un service est démarré.
\end{itemize}

\begin{remarque}
Il faut \textbf{redémarrer} (\emph{Restart}) Apache à chaque fois qu'on modifie un
fichier de configuration (comme \code{httpd.conf}), sinon le changement n'est pas pris
en compte. En revanche, modifier une simple page \code{.php} ne nécessite aucun
redémarrage : il suffit de recharger la page dans le navigateur.
\end{remarque}

\section{Le répertoire racine}

\begin{definition}
Le \textbf{répertoire racine} (\emph{document root}) est le dossier dans lequel le
serveur web va chercher les pages à servir. Tout fichier placé dans ce dossier
devient accessible depuis le navigateur via \code{http://localhost/}.
\end{definition}

Le nom du répertoire racine dépend de l'environnement :
\begin{itemize}
  \item \textbf{WampServer} : \code{www} (par exemple \code{C:\textbackslash wamp64\textbackslash www}) ;
  \item \textbf{XAMPP} : \code{htdocs} (par exemple \code{C:\textbackslash xampp\textbackslash htdocs}) ;
  \item \textbf{LAMP (Linux)} : \code{/var/www/html}.
\end{itemize}

\begin{exemple}
Sous WampServer, le fichier \code{C:\textbackslash wamp64\textbackslash www\textbackslash bulletin.php}
est accessible à l'adresse \code{http://localhost/bulletin.php}. Si on le range dans
un sous-dossier \code{lycee}, l'adresse devient \code{http://localhost/lycee/bulletin.php}.
\end{exemple}

\begin{attention}
Une page placée \textbf{en dehors} du répertoire racine (par exemple sur le Bureau)
ne pourra \textbf{pas} être servie par Apache : le navigateur ne la trouvera pas à
l'adresse \code{http://localhost/}. Place toujours tes fichiers \code{.php} dans
\code{www} (ou \code{htdocs}).
\end{attention}

\section{Le port d'écoute}

\begin{definition}
Un \textbf{port} est un numéro qui identifie un service réseau sur une machine.
Le serveur web Apache « écoute » par défaut sur le \textbf{port 80}, le port standard
du protocole HTTP. C'est pourquoi on peut écrire \code{http://localhost} sans préciser
de numéro.
\end{definition}

\begin{remarque}
Quand aucun port n'est précisé, le navigateur utilise automatiquement le port 80 pour
\code{http}. Écrire \code{http://localhost} revient donc exactement à écrire
\code{http://localhost:80}.
\end{remarque}

\subsection{Identifier et modifier le port}

Le port d'écoute d'Apache est défini dans son fichier de configuration principal,
\code{httpd.conf}. On y trouve une ligne \code{Listen} :

\begin{center}\small
\begin{tabular}{@{}ll@{}}
\toprule
\textbf{Environnement} & \textbf{Emplacement de \code{httpd.conf}}\\
\midrule
WampServer & \code{C:\textbackslash wamp64\textbackslash bin\textbackslash apache\textbackslash apache2.x.x\textbackslash conf}\\
XAMPP      & \code{C:\textbackslash xampp\textbackslash apache\textbackslash conf}\\
\bottomrule
\end{tabular}
\end{center}

Dans WampServer, on peut ouvrir directement ce fichier par le menu de l'icône :
\emph{Apache} $\rightarrow$ \emph{httpd.conf}. La ligne à repérer est :

\begin{verbatim}
# Fichier httpd.conf - port d'ecoute d'Apache
Listen 80
\end{verbatim}

\fig{La directive \code{Listen} fixe le port sur lequel Apache attend les requêtes.}

Pour faire écouter Apache sur un autre port (par exemple le \textbf{8080}, souvent
utilisé quand le port 80 est déjà occupé), on remplace la valeur :

\begin{verbatim}
# Apache ecoute desormais sur le port 8080
Listen 8080
\end{verbatim}

\begin{methode}
\textbf{Changer le port d'écoute d'Apache.}
\begin{enumerate}[label=\arabic*)]
  \item Ouvrir le fichier \code{httpd.conf} (menu \emph{Apache}).
  \item Remplacer la ligne \code{Listen 80} par \code{Listen 8080}.
  \item Si elle existe, modifier aussi la ligne \code{ServerName localhost:80} en
        \code{ServerName localhost:8080}.
  \item \textbf{Enregistrer} le fichier et \textbf{redémarrer} Apache.
  \item Accéder désormais au site avec le port dans l'URL :
        \code{http://localhost:8080}.
\end{enumerate}
\end{methode}

\begin{attention}
Après avoir changé le port, l'adresse \code{http://localhost} (port 80) ne fonctionne
plus : il faut impérativement préciser le nouveau port, comme dans
\code{http://localhost:8080}. Et il ne faut pas oublier de redémarrer Apache, sinon
l'ancien réglage reste actif.
\end{attention}

\section{Héberger un site}

\subsection{En local (sur sa propre machine)}

Héberger un site « en local », c'est le faire fonctionner sur son propre ordinateur,
sans connexion Internet. C'est la méthode idéale pour développer et tester avant la
mise en ligne.

\begin{methode}
\textbf{Héberger un site en local.}
\begin{enumerate}[label=\arabic*)]
  \item \textbf{Placer le dossier} du site dans le répertoire racine. Exemple : copier
        le dossier \code{notesonline} dans \code{C:\textbackslash wamp64\textbackslash www}.
  \item \textbf{Démarrer les services} : lancer Apache (serveur web) et MySQL (si le
        site utilise une base de données). L'icône WampServer doit être verte.
  \item \textbf{Ouvrir le navigateur} et saisir l'adresse du site :
        \code{http://localhost/notesonline/}. Le serveur cherche alors un fichier
        \code{index.php} (ou \code{index.html}) dans ce dossier.
  \item \textbf{Naviguer} dans le site comme s'il était en ligne, mais visible
        uniquement depuis cette machine.
\end{enumerate}
\end{methode}

\begin{remarque}
L'adresse \code{localhost} (ou son équivalent numérique \code{127.0.0.1}) désigne
toujours « ma propre machine ». Un site en local n'est donc pas accessible depuis
Internet : il faut un véritable hébergeur pour le rendre public.
\end{remarque}

\subsection{En ligne (mise en production)}

Une fois le site terminé et testé en local, on le publie sur Internet pour qu'il soit
accessible à tous. Il faut alors deux choses :

\begin{itemize}
  \item un \textbf{hébergeur} : une entreprise qui loue de l'espace sur un serveur web
        allumé en permanence et connecté à Internet (par exemple OVH, Ionos, ou un
        hébergeur local camerounais comme Camtel/Blue) ;
  \item un \textbf{nom de domaine} : l'adresse facile à retenir du site (par exemple
        \code{www.onbuch.cm}), à louer auprès d'un bureau d'enregistrement. Le suffixe
        \code{.cm} est réservé aux sites du Cameroun.
\end{itemize}

\begin{exemple}
Un lycée de Yaoundé développe son site de publication des résultats. Il le teste
d'abord en local sous WampServer (\code{http://localhost/resultats/}). Une fois prêt,
il loue le nom de domaine \code{lycee-bilingue.cm} et un hébergement, puis transfère
ses fichiers \code{.php} sur le serveur de l'hébergeur (généralement par \textbf{FTP}).
Le site est alors accessible depuis n'importe quel ordinateur du monde.
\end{exemple}

\begin{remarque}
Pour transférer les fichiers vers l'hébergeur, on utilise le plus souvent le protocole
\textbf{FTP} (\emph{File Transfer Protocol}) avec un logiciel comme FileZilla, ou bien
l'interface d'administration fournie par l'hébergeur.
\end{remarque}

\section{Exercices}

\rubrique{Application}

\exo{Client ou serveur ?}\diff{1}
Pour chaque élément, indique s'il joue le rôle de \textbf{client} ou de
\textbf{serveur} : a) Mozilla Firefox ; b) Apache ; c) le navigateur d'un téléphone ;
d) Nginx ; e) Google Chrome.

\exo{Vrai ou faux ?}\diff{1}
Réponds par vrai ou faux et justifie : a) le navigateur exécute le code PHP ;
b) Apache est un serveur web ; c) HTTP est le protocole utilisé entre le navigateur
et le serveur web ; d) le client envoie la réponse et le serveur la requête ;
e) un site en local est visible depuis Internet.

\exo{Associer composant et rôle}\diff{1}
Associe chaque composant d'un environnement WAMP à son rôle :
\begin{enumerate}[label=\arabic*)]
  \item Apache \qquad a) exécute le code et fabrique le HTML
  \item MySQL  \qquad b) stocke et gère les données
  \item PHP    \qquad c) reçoit les requêtes et renvoie les pages
\end{enumerate}

\exo{Le bon répertoire racine}\diff{1}
Donne le nom du répertoire racine pour : a) WampServer ; b) XAMPP ; c) LAMP sous
Linux. Où faut-il placer un fichier \code{accueil.php} sous WampServer pour qu'il soit
accessible à \code{http://localhost/accueil.php} ?

\exo{Décoder une URL}\diff{2}
On donne l'adresse \code{http://localhost:8080/ecole/notes.php}. Indique :
a) le port d'écoute utilisé ; b) le sous-dossier du site dans le répertoire racine ;
c) le nom du fichier demandé ; d) pourquoi le port \code{8080} apparaît ici alors qu'il
n'apparaît pas dans \code{http://localhost/}.

\rubrique{Entraînement}

\exo{Changer le port d'écoute}\diff{2}
Le port 80 est déjà occupé par un autre logiciel et Apache refuse de démarrer (icône
orange). Décris la procédure complète pour faire écouter Apache sur le port \code{8081},
en précisant le fichier modifié, la ligne à changer et l'adresse à utiliser ensuite.

\exo{Procédure d'hébergement local}\diff{2}
Un camarade t'envoie un dossier \code{cantine} contenant son site PHP. Décris, étape
par étape, ce que tu dois faire pour l'exécuter sur ta machine avec WampServer, depuis
la copie du dossier jusqu'à l'affichage dans le navigateur.

\exo{Le cycle requête/réponse}\diff{2}
Un élève saisit \code{http://www.onbuch.cm/resultats.php} dans son navigateur. Décris,
dans l'ordre, les quatre étapes du cycle requête/réponse HTTP jusqu'à l'affichage de la
page, en précisant à quel moment le code PHP est exécuté.

\exo{De local à en ligne}\diff{3}
Le club informatique d'un lycée de Douala a terminé son site sous XAMPP
(\code{http://localhost/club/}). a) Pourquoi ce site n'est-il pas encore visible par
les élèves chez eux ? b) Cite les \textbf{deux} éléments à acquérir pour le mettre en
ligne. c) Par quel moyen transfère-t-on les fichiers vers l'hébergeur ?

\section{Corrigés}

\corrige{de l'exercice 1}
Le \textbf{client} est le logiciel qui demande et affiche les pages (un navigateur) ;
le \textbf{serveur} est le logiciel qui héberge et renvoie les pages.
\begin{itemize}
  \item a) Mozilla Firefox $\rightarrow$ \textbf{client} (navigateur) ;
  \item b) Apache $\rightarrow$ \textbf{serveur} (serveur web) ;
  \item c) navigateur d'un téléphone $\rightarrow$ \textbf{client} ;
  \item d) Nginx $\rightarrow$ \textbf{serveur} (serveur web) ;
  \item e) Google Chrome $\rightarrow$ \textbf{client} (navigateur).
\end{itemize}

\corrige{de l'exercice 2}
\begin{enumerate}[label=\alph*)]
  \item \textbf{Faux} : le PHP s'exécute sur le \emph{serveur}, pas dans le navigateur.
  \item \textbf{Vrai} : Apache est bien un serveur web.
  \item \textbf{Vrai} : HTTP est le protocole d'échange entre client et serveur web.
  \item \textbf{Faux} : c'est le client qui envoie la \emph{requête} et le serveur qui
        renvoie la \emph{réponse}.
  \item \textbf{Faux} : un site en local (\code{localhost}) n'est visible que sur la
        machine elle-même, pas depuis Internet.
\end{enumerate}

\corrige{de l'exercice 3}
\begin{itemize}
  \item 1 $\rightarrow$ c) Apache reçoit les requêtes et renvoie les pages ;
  \item 2 $\rightarrow$ b) MySQL stocke et gère les données ;
  \item 3 $\rightarrow$ a) PHP exécute le code et fabrique le HTML.
\end{itemize}

\corrige{de l'exercice 4}
Le répertoire racine est : a) \code{www} pour WampServer ; b) \code{htdocs} pour
XAMPP ; c) \code{/var/www/html} pour LAMP sous Linux. Sous WampServer, le fichier
\code{accueil.php} doit être placé directement dans le dossier \code{www}
(c'est-à-dire \code{C:\textbackslash wamp64\textbackslash www\textbackslash accueil.php})
pour être accessible à \code{http://localhost/accueil.php}.

\corrige{de l'exercice 5}
Dans \code{http://localhost:8080/ecole/notes.php} :
\begin{enumerate}[label=\alph*)]
  \item le port d'écoute est \textbf{8080} (la valeur après les deux-points) ;
  \item le sous-dossier du site est \code{ecole} (placé dans le répertoire racine) ;
  \item le fichier demandé est \code{notes.php} ;
  \item le port \code{8080} apparaît parce qu'Apache a été configuré sur ce port (au
        lieu du 80 par défaut). Dans \code{http://localhost/}, le port n'est pas écrit
        car le navigateur utilise automatiquement le port 80, valeur par défaut de HTTP.
\end{enumerate}

\corrige{de l'exercice 6}
Procédure pour faire écouter Apache sur le port 8081 :
\begin{enumerate}[label=\arabic*)]
  \item ouvrir le fichier de configuration \code{httpd.conf} d'Apache (menu
        \emph{Apache} $\rightarrow$ \emph{httpd.conf} dans WampServer) ;
  \item remplacer la ligne \code{Listen 80} par \code{Listen 8081} (et, si présente,
        \code{ServerName localhost:80} par \code{ServerName localhost:8081}) ;
  \item enregistrer le fichier ;
  \item redémarrer Apache (\emph{Restart Services}) pour appliquer le changement ;
  \item accéder au site avec le nouveau port : \code{http://localhost:8081/}.
\end{enumerate}

\corrige{de l'exercice 7}
Pour exécuter le dossier \code{cantine} sous WampServer :
\begin{enumerate}[label=\arabic*)]
  \item copier le dossier \code{cantine} dans le répertoire racine
        \code{C:\textbackslash wamp64\textbackslash www} ;
  \item démarrer WampServer et lancer les services Apache (et MySQL si le site utilise
        une base de données) : l'icône doit devenir verte ;
  \item ouvrir le navigateur et saisir \code{http://localhost/cantine/} ;
  \item le serveur affiche le fichier \code{index.php} du dossier ; on peut alors
        naviguer dans le site.
\end{enumerate}

\corrige{de l'exercice 8}
Cycle requête/réponse pour \code{http://www.onbuch.cm/resultats.php} :
\begin{enumerate}[label=\arabic*)]
  \item \textbf{Requête} : le navigateur (client) envoie une requête HTTP au serveur
        pour demander la page \code{resultats.php} ;
  \item \textbf{Traitement} : le serveur web reçoit la requête et fait \textbf{exécuter
        le code PHP} de la page (c'est à ce moment que le PHP s'exécute, éventuellement
        en interrogeant une base de données) ;
  \item \textbf{Réponse} : le serveur renvoie au navigateur la page \textbf{HTML}
        produite par le PHP ;
  \item \textbf{Affichage} : le navigateur interprète le HTML reçu et affiche la page à
        l'élève.
\end{enumerate}

\corrige{de l'exercice 9}
\begin{enumerate}[label=\alph*)]
  \item Le site n'est visible que sur la machine du club car \code{localhost} désigne
        cette seule machine ; il n'est pas accessible depuis Internet.
  \item Pour le mettre en ligne, il faut acquérir : un \textbf{hébergeur} (un serveur
        connecté à Internet en permanence) et un \textbf{nom de domaine}
        (par exemple \code{club-douala.cm}).
  \item On transfère les fichiers vers l'hébergeur par \textbf{FTP} (avec un logiciel
        comme FileZilla) ou via l'interface d'administration de l'hébergeur.
\end{enumerate}

\begin{aretenir}
Sur le Web, le \textbf{client} (navigateur) envoie une \textbf{requête HTTP} et le
\textbf{serveur web} (Apache) renvoie une \textbf{réponse} : la page HTML. Un
environnement intégré (WampServer, XAMPP, LAMP) réunit \textbf{Apache + MySQL + PHP}.
On place ses pages dans le \textbf{répertoire racine} (\code{www} ou \code{htdocs}),
Apache écoute par défaut sur le \textbf{port 80} (modifiable via \code{Listen} dans
\code{httpd.conf}), et on teste tout en local sur \code{http://localhost/} avant la
mise en ligne (hébergeur + nom de domaine).
\end{aretenir}
